\documentclass[12pt]{article}
\usepackage[T1]{fontenc}
\usepackage[utf8]{inputenc}
\usepackage{lmodern}
\usepackage{textcomp}
\usepackage{enumitem}
\usepackage{geometry}
\geometry{margin=1in}
\begin{document}
\begin{center}
Answer Key - Constitutional Law - Undergraduate Level Assessment 11 \
\small (Answers are ordered exactly as in the question paper.)
\end{center}

\section*{Multiple Choice}
\begin{enumerate}[label=\arabic*.]
\item \textbf{Answer:} A
\\ \textit{Options:} A) Historical; B) Analytical; C) Command; D) Sociological
\\ \textit{Note:} The Historical School of jurisprudence posits that law evolves from the customs and traditions of a society, emphasizing the importance of historical context in understanding legal principles.
\item \textbf{Answer:} C
\\ \textit{Options:} A) Command, sovereign and enforceability; B) Command, sovereign and legal remedy; C) Command, sovereign and sanction; D) Command, sovereign and obedience by subject
\\ \textit{Note:} The correct answer identifies the key attributes of law as defined by legal positivist John Austin, who emphasized that law is a command issued by a sovereign backed by sanctions, thereby framing the nature and authority of legal systems.
\item \textbf{Answer:} D
\\ \textit{Options:} A) They adopt different attitudes towards the role of the courts in maintaining order.; B) They disagree about the role of law in society.; C) They have opposing views about the nature of contractual obligations.; D) They differ in respect of their account of life before the social contract.
\\ \textit{Note:} The correct answer highlights that Hobbes views the state of nature as a chaotic and violent environment requiring a strong sovereign for order, while Locke perceives it as a generally peaceful state that can lead to the establishment of a government to protect natural rights, illustrating their fundamentally different views on human nature and society prior to the social contract.
\item \textbf{Answer:} C
\\ \textit{Options:} A) That there is no distinction between law and morality.; B) That there is a distinction between right and wrong.; C) That facts about the world or human nature cannot normally ordain what ought to be; D) That human rights are fundamentally unsound.
\\ \textit{Note:} Dworkin's endorsement of Hume's principle emphasizes the distinction between descriptive facts and prescriptive norms, arguing that moral and ethical judgments cannot be solely derived from empirical observations of the world or human nature.
\item \textbf{Answer:} D
\\ \textit{Options:} A) Authority.; B) Charisma.; C) Co-operation.; D) Capitalism.
\\ \textit{Note:} Weber's analysis links the development of formally rational law to the rise of capitalism, arguing that the economic imperatives of capitalist society necessitated a legal system characterized by predictability, calculability, and efficiency.
\end{enumerate}

\section*{True / False}
\begin{enumerate}[label=\arabic*.]
\setcounter{enumi}{5}
\item \textbf{Answer:} False
\\ \textit{Options:} A) True; B) False
\\ \textit{Note:} Unlawful homicide, regardless of the perpetrator's status as a minister or the location of the act, does not qualify as an official act (jure imperii) and therefore does not attract immunity from prosecution under international law.
\item \textbf{Answer:} False
\\ \textit{Options:} A) True; B) False
\\ \textit{Note:} A contract for the sale of illegal substances is unenforceable because it violates public policy and is considered void ab initio, regardless of mutual consent between the parties.
\item \textbf{Answer:} True
\\ \textit{Options:} A) True; B) False
\\ \textit{Note:} Collective security is founded on the principle that member states of an organization, like the United Nations, agree to respond collectively to threats to peace, which includes the authorization of armed force by the UN Security Council when necessary to maintain international peace and security.
\item \textbf{Answer:} True
\\ \textit{Options:} A) True; B) False
\\ \textit{Note:} The case of Goodwin v UK (2002) illustrates the 'living instrument principle' as the European Court of Human Rights interpreted the right to respect for private life in a way that reflected contemporary societal understandings of gender identity and rights, thereby adapting to evolving norms.
\item \textbf{Answer:} True
\\ \textit{Options:} A) True; B) False
\\ \textit{Note:} Standard-setting in human rights diplomacy is characterized by the establishment of non-binding legal norms that serve as guidelines for states and organizations to promote and uphold human rights standards without imposing legal obligations.
\end{enumerate}

\section*{Long Form Response}
\begin{enumerate}[label=\arabic*.]
\setcounter{enumi}{10}
\item \textbf{Answer:} Ronald Dworkin critiques Richard Posner's economic analysis of law by challenging the notion that wealth maximization should be considered a fundamental value in legal reasoning. Dworkin argues that reducing legal principles to mere economic efficiency overlooks the moral and ethical dimensions inherent in the law. He contends that wealth, while significant in certain contexts, should not be the ultimate goal of legal frameworks, as it fails to account for justice, individual rights, and the intrinsic value of human dignity. By prioritizing wealth, Posner's approach risks undermining the broader principles of fairness and equality that are essential in a constitutional democracy. Thus, Dworkin asserts that law must reflect more than just economic calculations; it should embody values that uphold the integrity of society as a whole.
\\ \textit{Note:} correct because it accurately captures Dworkin's argument that prioritizing wealth maximization in legal analysis neglects essential moral and ethical considerations, which are crucial for achieving justice and protecting individual rights within the legal framework.
\item \textbf{Answer:} A state can express its consent to be bound by a treaty through several methods, including signature, ratification, acceptance, approval, or accession. Each method has distinct implications within constitutional law: a signature may indicate an intention to be bound but typically requires further action to achieve full consent; ratification is the formal approval process often requiring legislative endorsement; acceptance and approval can signify a state's formal agreement to the treaty terms, while accession allows a state to join an existing treaty without negotiation. The choice of method can affect the treaty's domestic implementation and the extent to which it influences national law, as some methods may necessitate additional legislative action, thereby highlighting the interplay between international obligations and constitutional frameworks. Overall, the method of consent reflects the state's commitment to uphold international agreements within its legal system.
\\ \textit{Note:} correct because it accurately identifies the various methods a state can use to consent to a treaty, such as signature, ratification, acceptance, approval, and accession, while also highlighting the constitutional implications of each method, particularly the varying degrees of commitment and legislative involvement required for full consent.
\end{enumerate}

\vfill
\noindent\textit{End of Answer Key}
\end{document}